%! TeX program = lualatex
%! TEX options = -synctex=1 -interaction=nonstopmode -file-line-error --shell-escape "%DOC%"
\documentclass[10pt,a5paper]{mszalik}

\begin{document}

\thispagestyle{empty}

\begin{center}
	\vspace*{0.5cm}


	\bfseries\scshape
	\huge ŚRODA POPIELCOWA

	\vfill

	\pgfornament[width=7cm]{88}

	\vfill

	\begin{figure}[h]
		\centering
		\includegraphics[width=0.5\linewidth]{logo.png}
	\end{figure}

	% \vfill

	{\large Duszpasterstwo Wiernych Tradycji Łacińskiej \\ w Archidiecezji Wrocławskiej}

	\bigskip

	{\Large \textbf{2020}}


\end{center}

\afterpage{\null\thispagestyle{empty}\newpage}

% \newpage
% \thispagestyle{empty}

\vspace*{2cm}
"Popiół jest symbolem znikomości świata oraz pokory i pokuty. Poddając się
obrzędowi posypania głów popiołem, uznajemy publicznie naszą grzeszność i wyrażamy
chęć poprawy. Na mocy zasług Kościoła otrzymujemy pomoc Bożą do gorliwego odbycia
Wielkiego Postu"

\bigskip

\hfill \textit{Mszał Rzymski, Wydawnictwo Pallotinum, Poznań, 1963}

\vspace*{\fill}

\hrulefill

{\footnotesize Niniejszy mszalik przygotowany jest w celu wyjaśnienia obrzędów i
	edukacji liturgicznej wiernych. Teksty liturgiczne wraz z tłumaczeniami oraz
	grafiki pochodzą z:

	\begin{itemize}[leftmargin=*]
		\item \textit{Mszału Rzymskiego z dodaniem nabożeństw nieszpornych}
		      autorstwa O. Gaspara Lefebvre (Poznań 1932)
		\item \textit{Mszału Rzymskiego} autorstwa benedyktynów tynieckich
		      (Poznań 1963)
		\item \textit{Missale Romanum} wyd. Typis Polyglotis Vaticanis (1927)
		      oraz Pustet (1899)
		\item \textit{Cantus Selecti} (1957)
		\item \textit{Liber Usualis} (1961)
	\end{itemize}

	W niektórych przypadkach wykorzystano cyfrową wersję powyższych tekstów
	opublikowaną na portalu \texttt{divinumofficium.com}.

	\medskip

	Cytaty z homilii papieża Benedykta XVI za \texttt{opoka.org.pl}.

	\medskip

	Teksty czytań opatrzone komentarzem i wyjaśnieniem zaczerpnięto z
	\textit{Biblii Tysiąclecia Online} (Poznań 2003). Zostały one wykorzystane
	na mocy prawa cytatu. (art. 29. Ustawy o prawie autorskim i prawach
	pokrewnych) Prawa autorskie do nich posiada wyd. Pallotinum. Tekst Biblii
	Tysiąclecia jest oficjalnym tekstem Kościoła w Polsce zatwierdzonym do
	stosowania w liturgii.

	\bigskip

	opracowanie: Piotr Matuszak\\
	nadzór merytoryczny ks. dr Ireneusz Bakalarczyk\\
	skład i łamanie: Michał Siemaszko\\

	email: tradycja@archidiecezja.wroc.pl
}
\hrule

%	\afterpage{\null\thispagestyle{empty}\newpage}
\newpage

%%%%%%%%%%%%%%%%%%%%%%%%%%%%%%%%%%%%%%%%%%%%%%%%%%%%%%%%%%%%%%%%%%%%%%%%%%

\newpage

\fancyhead[C]{POŚWIĘCENIE POPIOŁU}

\begin{figure}[h]
	\centering
	\includegraphics[width=\linewidth]{Figures/popielec.jpg}
\end{figure}

\section{POŚWIĘCENIE POPIOŁU}

% \rubric{Krótki opis co się będzie działo}

\proroctwo{Pieśń na wejście}

\gregorioscore{Melody/PD.gabc}

\newpage

\oremuss{
2. N\'{o}stris m\'{a}lis off\'{e}ndimus \\
T\'{u}am D\'{e}us clem\'{e}ntiam \\
Eff\'{u}nde n\'{o}bis d\'{e}super \\
Rem\'{i}ssor indulg\'{e}ntiam. \rrit \textit{P\'{a}rce D\'{o}mine.}

\medskip

3. Dans t\'{e}mpus accept\'{a}bile, \\
Da lacrim\'{a}rum r\'{i}vulis \\
Lav\'{a}re c\'{o}rdis v\'{i}ctimam, \\
Quam la\'{e}ta ad\'{u}rat c\'{a}ritas. \rrit \textit{P\'{a}rce D\'{o}mine.}
}{
4. Audi, ben\'{i}gne C\'{o}nditor, \\
N\'{o}stras pr\'{e}ces cum fl\'{e}tibus \\
In hoc s\'{a}cro jej\'{u}nio \\
F\'{u}sas quadragen\'{a}rio. \rrit \textit{P\'{a}rce D\'{o}mine.}

\medskip

5. Scrut\'{a}tor \'{a}lme c\'{o}rdium, \\
Infirma tu scis v\'{i}rium; \\
Ad te rev\'{e}rsis \'{e}xhibe \\
Remissi\'{o}nis gr\'{a}tiam. \rrit \textit{P\'{a}rce D\'{o}mine.}
}


\oremus{Antyfona (Ps 68:17.2)}{Exáudi nos, Dómine, quóniam benígna est misericórdia
	tua: secúndum multitúdinem miseratiónum tuárum réspice nos, Dómine.\\
	Salvum me fac, Deus: quóniam intravérunt aquæ usque ad ánimam meam. \medskip

	\vv Glória Patri, et Fílio, et Spirítui Sancto. \\\\
	\rr Sicut erat in princípio, et nunc, et semper, et in sǽcula sæculórum. Amen.}
{Wysłuchaj nas, Panie, gdyż łaskawe jest miłosierdzie Twoje, w wielkiej
	litości swojej wejrzyj na nas.\\

	Wybaw mnie, Boże, bo wody doszły mi po szyję. \medskip

	\vv Chwała Ojcu, i Synowi i Duchowi Świętemu.\\
	\rr Jak była na początku, teraz i zawsze i na wieki wieków. Amen.
}


\oremus{Antyfona 1. (Jl 2:13)}{Immutémur habitu, in cínere et cilício: ieiunémus,
	et plorémus ante Dóminum: quia multum miséricors est dimíttere peccáta nostra
	Deus noster.}{Odmieńmy sposób życia, pokutujmy w popiele i włosiennicy, pośćmy
	i płaczmy przed Panem, gdyż Bóg nasz jest pełen miłosierdzia i gotów odpuścić
	nam nasze grzechy.}

\oremus{Antyfona 2. (Jl 2:17, Est 13:17)}{Inter vestíbulum et altáre plorábunt
	sacerdótes minístri Dómini, et dicent: Parce, Dómine, parce pópulo tuo:
	et ne claudas ora canéntium te, Dómine.}{Między przedsionkiem a ołtarzem
	niechaj płaczą kapłani, słudzy Pańscy, i niech mówią: Przepuść, Panie, przepuść
	ludowi Twojemu, a nie zamykaj ust śpiewającym chwałę Twoją, Panie.}

\oremus{Responsorium (Est 13, Jl 2, Ps 78:9)}{Emendémus in mélius, quæ ignoránter
	peccávimus: ne, subito præoccupáti die mortis, quærámus spátium pæniténtiæ,
	et inveníre non póssimus.\\

	\rr Atténde, Dómine, et miserére: quia peccávimus tibi.

	\vv Adiuva nos, Deus, salutáris noster: et propter honórem nóminis tui,
	Dómine, líbera nos.\\
	\rr Atténde, Dómine, et miserére: quia peccávimus tibi.

	\vv Glória Patri, et Fílio, et Spirítui Sancto.\\\\
	\rr Atténde, Dómine, et miserére: quia peccávimus tibi.
}{
	Zmieńmy na lepsze wszystko to, czym zgrzeszyliśmy w naszym zaślepieniu,
	abyśmy, gdy nas znienacka zaskoczy dzień śmierci, nie szukali czasu
	na pokutę, kiedy go już nie będziemy mogli znaleźć.

	\rr Wejrzyj na nas, Panie, i zmiłuj się: albowiem zgrzeszyliśmy przeciw Tobie.

	\vv Wspomóż nas Boże, nasz Zbawicielu, i dla chwały imienia Twego wybaw nas, Panie.	\\
	\rr Wejrzyj na nas, Panie, i zmiłuj się: albowiem zgrzeszyliśmy przeciw Tobie.

	\vv Chwała Ojcu, i Synowi, i Duchowi Świętemu.	\\
	\rr Wejrzyj na nas, Panie, i zmiłuj się: albowiem zgrzeszyliśmy przeciw Tobie.}

\oremus{Pieśń}{
	1. Posypmy głowy popiołem, \\
	uderzmy przed Panem czołem; \\
	Zapustne śmiechy na stronę, \\
	cierniową wijmy koronę. \\

	\medskip

	2. Posypmy głowy popiołem, \\
	grzmi niebo głosem surowym: \\
	"Pokutę czyńcie za grzechy, \\
	na stronę teraz uciechy". \\

	\medskip

	3. Posypmy głowy popiołem, \\
	otoczmy Pana pospołem, \\
	Bo Pański sąd sprawiedliwy, \\
	a dla grzesznika straszliwy.

}{
	4. Posypmy głowy popiołem, \\
	już Zbawca cierpieć gotowy, \\
	Ponosić męki, katusze, \\
	by nasze odkupić dusze.\\

	\medskip

	5. Posypmy głowy popiołem, \\
	o śmierci myślmy pospołem; \\
	Nim głosu Twego wezwanie \\
	na sąd Twój stawi nas, Panie.\\

	\medskip

	6. Posypmy głowy popiołem, \\
	z Chrystusem nieśmy krzyż społem, \\
	Abyśmy z Nim zmartwychwstali \\
	i hymn Mu chwały śpiewali
}

\oremus{Modlitwa}{
	\vv Dóminus vobíscum.\\
	\rr Et cum spíritu tuo.

	Orémus. Concéde nobis, Dómine, præsídia milítiæ christiánæ sanctis inchoáre ieiúniis:
	ut, contra spiritáles nequítias pugnatúri, continéntiæ muniámur auxíliis.
	Per Christum, Dóminum nostrum.

	\rr Amen.}{
	\vv Pan z wami.\\
	\rr I z duchem twoim.

	Módlmy się. Dozwól nam, Panie, przez post zająć stanowisko w chrześcijańskim boju,
	abyśmy zostali umocnieni wstrzemięźliwością do walki ze złym duchem.
	Przez Chrystusa, Pana naszego.\\

	\rr Amen.}

\fancyhead[C]{MSZA ŚWIĘTA}

\section{ŚRODA POPIELCOWA}
 {\station{Stacja u św. Sabiny}}



\oremus{Introit (Sap 11:24 11:25; 11:27, Ps 56:2)}{ Miseréris ómnium, Dómine, et nihil
	odísti eórum quæ fecísti, dissímulans peccáta hóminum propter pæniténtiam et
	parcens illis: quia tu es Dóminus, Deus noster\\

	Miserére mei, Deus, miserére mei: quóniam in te confídit ánima mea. \medskip

	\vv Glória Patri, et Fílio, et Spirítui Sancto. \\\\
	\rr Sicut erat in princípio, et nunc, et semper, et in sǽcula sæculórum. Amen.}{
	Wszystkim okazujesz miłosierdzie, Panie, i żadnego ze swych stworzeń nie masz w
	nienawiści. Darowujesz grzechy ludziom, aby się nawrócili, i przebaczasz im:
	albowiem Ty jesteś Panem, Bogiem naszym.

	Zmiłuj się nade mną, Boże, zmiłuj się nade mną, bo Tobie ufa dusza moja.\medskip

	\vv Chwała Ojcu, i Synowi i Duchowi Świętemu.\\
	\rr Jak była na początku, teraz i zawsze i na wieki wieków. Amen.
}

\oremus{Kolekta}{Præsta, Dómine, fidélibus tuis: ut ieiuniórum veneránda
	sollémnia, et cóngrua pietáte suscípiant, et secúra devotióne
	percúrrant.}{Udziel, Panie, wiernym Twoim łaski, aby rozpoczęli ten czcigodny i
	uroczysty post z należną pobożnością i przebyli go służąc Tobie w pokoju.}

\proroctwo{{Lekcja (Jl 2:12-19 )}}

\rubric{Kościół przypomina nam przede wszystkim mocne wezwanie, jakie prorok Joel
	skierował do narodu Izraelskiego (...). Podkreślić należy wyrażenie «całym swym sercem»,
	które oznacza, że nawrócenie pochodzi z głębi naszych myśli i uczuć, leży u podstaw
	naszych decyzji, wyborów i działań, w geście całkowitej i radykalnej wolności.
	Ale czy możliwy jest ten powrót do Boga? Tak, jest możliwy, bo istnieje siła,
	która nie pochodzi z naszego serca, ale która uwalnia się w sercu samego Boga.
	Jest to siła Jego miłosierdzia. --- Benedykt XVI, 13.02.2013}

\oremuss{Hæc dicit Dóminus: Convertímini ad me in toto corde vestro, in ieiúnio,
	et in fletu, et in planctu. Et scíndite corda vestra, et non vestiménta vestra,
	et convertímini ad Dóminum, Deum vestrum: quia benígnus et miséricors est,
	pátiens, et multæ misericórdiæ, et præstábilis super malítia. Quis scit, si
	convertátur, et ignóscat, et relínquat post se benedictiónem, sacrifícium et
	libámen Dómino, Deo vestro? Cánite tuba in Sion, sanctificáte ieiúnium, vocáte
	coetum, congregáte pópulum, sanctificáte ecclésiam, coadunáte senes, congregáte
	parvulos et sugéntes úbera: egrediátur sponsus de cubíli suo, et sponsa de
	thálamo suo. Inter vestíbulum et altare plorábunt sacerdótes minístri Dómini, et
	dicent: Parce, Dómine, parce pópulo tuo: et ne des hereditátem tuam in
	oppróbrium, ut dominéntur eis natiónes. Quare dicunt in pópulis: Ubi est Deus
	eórum? Zelátus est Dóminus terram suam, et pepércit pópulo suo. Et respóndit
	Dóminus, et dixit populo suo: Ecce, ego mittam vobis fruméntum et vinum et
	óleum, et replebímini eis: et non dabo vos ultra oppróbrium in géntibus: dicit
	Dóminus omnípotens.}{
	To mówi Pan: «Nawróćcie się do Mnie całym swym sercem, przez post i
	płacz, i lament». Rozdzierajcie jednak serca wasze, a nie szaty! Nawróćcie się
	do Pana Boga waszego! On bowiem jest łaskawy, miłosierny, nieskory do
	gniewu i wielki w łaskawości, a lituje się na widok niedoli. Kto wie? Może
	znów pożałuje i pozostawi po sobie błogosławieństwo plonów na ofiarę
	z pokarmów i ofiarę płynną dla Pana Boga waszego.
	Na Syjonie dmijcie w róg, zarządźcie święty post, ogłoście uroczyste
	zgromadzenie. Zbierzcie lud, zwołajcie świętą społeczność, zgromadźcie
	starców, zbierzcie dzieci, i ssących piersi! Niech wyjdzie oblubieniec ze swojej
	komnaty a oblubienica ze swego pokoju! Między przedsionkiem
	a ołtarzem niechaj płaczą kapłani, słudzy Pańscy! Niech mówią: «Przepuść,
	Panie, ludowi Twojemu i nie daj dziedzictwa swego na pohańbienie, aby
	poganie nie zapanowali nad nami. Czemuż mówić mają między narodami:
	Gdzież jest ich Bóg?»
	I Pan zapalił się zazdrosną miłością ku swojej ziemi, i zmiłował się nad
	swoim ludem. Odpowiedział Pan i rzekł swojemu ludowi: «Oto Ja wam posyłam
	zboże i moszcz, i oliwę, i nasyćcie się nimi, i nie wydam was więcej na
	pohańbienie między poganami.}

\proroctwo{Graduał (Ps 56:2; 56:4)}

\oremuss{Miserére mei, Deus, miserére mei: quóniam in te confídit ánima mea.

	\vv Misit de coelo, et liberávit me, dedit in oppróbrium conculcántes me.}
{Zmiłuj się nade mną, Boże, zmiłuj się nade mną, bo Tobie ufa dusza moja.

	\vv Sięgnął nieba i wybawił mnie, hańbą okrył tych, co mnie dręczą.}

\proroctwo{Traktus (Ps 102:10; 78:8-9)}

\oremuss{Dómine, non secúndum peccáta nostra, quæ fécimus nos: neque secúndum
iniquitátes nostras retríbuas nobis

Dómine, ne memíneris iniquitátum nostrarum antiquarum: cito antícipent nos
misericórdiæ tuæ, quia páuperes facti sumus nimis. {\footnotesize{Hic
	genuflectitur}} \medskip

\vv Adiuva nos, Deus, salutáris noster: et propter glóriam nóminis tui,
Dómine, libera nos: et propítius esto peccátis nostris, propter nomen
tuum.}{Panie, nie postępuj z nami według naszych grzechów ani nie odpłacaj
nam według win naszych.

Nie pamiętaj, Panie, dawnych nieprawości naszych. Niech rychło wyjdzie ku
nam miłosierdzie Twoje, bo jesteśmy bardzo nieszczęśliwi. {\footnotesize Tu
się przyklęka} \medskip

\vv Wspomóż nas, Boże, nasz Zbawicielu, i dla chwały imienia Twego wybaw
nas, Panie, i odpuść nam grzechy dla imienia Twego.}

\newpage

\proroctwo{Ewangelia ( Mt 6:16-21)}

\rubric{Jezus odnosi się do trzech podstawowych praktyk Prawa Mojżeszowego:
	jałmużna, modlitwa i post; są one również tradycyjnymi wskazówkami dla
	wielkopostnego wędrowania, by odpowiedzieć na zaproszenie do «powrotu do
	Boga całym sercem». Jezus jednak podkreśla, że tym, co decyduje o
	autentyczności wszelkich gestów religijnych jest jakość i prawdziwość
	relacji z Bogiem. Dlatego właśnie Jezus demaskuje religijną hipokryzję
	jako zachowanie, które promuje zewnętrzność, przyjmując postawy,
	które zabiegają o poklask i uznanie. --- Benedykt XVI, 13.02.2013}

\oremuss{In illo témpore: Dixit Iesus discípulis suis: Cum ieiunátis, nolíte
	fíeri, sicut hypócritæ, tristes. Extérminant enim fácies suas, ut appáreant
	homínibus ieiunántes. Amen, dico vobis, quia recepérunt mercédem suam. Tu
	autem, cum ieiúnas, unge caput tuum, et fáciem tuam lava, ne videáris
	homínibus ieiúnans, sed Patri tuo, qui est in abscóndito: et Pater tuus, qui
	videt in abscóndito, reddet tibi. Nolíte thesaurizáre vobis thesáuros in
	terra: ubi ærúgo et tínea demólitur: et ubi fures effódiunt et furántur.
	Thesaurizáte autem vobis thesáuros in coelo: ubi neque ærúgo neque tínea
	demólitur; et ubi fures non effódiunt nec furántur. Ubi enim est thesáurus
	tuus, ibi est et cor tuum.}{W owym czasie: Jezus powiedział do swoich uczniów: Kiedy pościcie,
	nie bądźcie posępni jak obłudnicy. Przybierają oni wygląd ponury, aby pokazać
	ludziom, że poszczą. Zaprawdę, powiadam wam: już odebrali swoją nagrodę.
	Ty zaś, gdy pościsz, namaść sobie głowę i umyj twarz, aby nie ludziom
	pokazać, że pościsz, ale Ojcu twemu, który jest w ukryciu. A Ojciec twój, który
	widzi w ukryciu, odda tobie.
	Nie gromadźcie sobie skarbów na ziemi, gdzie mól i rdza niszczą
	i gdzie złodzieje włamują się, i kradną. Gromadźcie sobie skarby w niebie, gdzie
	ani mól, ani rdza nie niszczą i gdzie złodzieje nie włamują się, i nie kradną. Bo
	gdzie jest twój skarb, tam będzie i serce twoje.}

\oremus{Offertorium (Ps 29:2-3)}{Exaltábo te, Dómine, quóniam suscepísti me, nec
	delectásti inimícos meos super me: Dómine, clamávi ad te, et sanásti
	me.}{Sławić Cię będę, Panie, bo mnie wybawiłeś i nie sprawiłeś ze mnie
	uciechy mym wrogom. Panie, do Ciebie wołałem, a Tyś mnie uzdrowił.}

\oremus{Pieśń}{
	1. Bądź mi litościw, Boże nieskończony,\\
	Według wielkiego miłosierdzia Twego.\\
	Według litości Twej niepoliczonéj\\
	Chciej zmazać mnóstwo przewinienia mego.

}{

	2. Obmyj mię z złości, obmyj tej godziny,\\
	Oczyść mię z brudu, w którym grzech mię trzyma;\\
	Bo ja poznaję wielkość mojej winy \\
	I grzech mój zawsze przed mymi oczyma. \\
}

\newpage

\oremuss{

	3. Odpuść! przed Tobą grzech mój popełniony, \\
	Boś przyrzekł, że ta kary ujdzie głowa, \\
	Którąć przyniesie grzesznik uniżony,\\
	By nie mówiono, że nie trzymasz słowa.

}{

	4. Odwróć twarz Twoją od przestępstwa mego\\
	I wszystkie moje pomaż nieprawości; \\
	Stwórz serce czyste, warte Boga swego, \\
	A ducha prawdy w moje wlej wnętrzności.\\
}
%
\begin{center}
	5. Nie oddalaj mię od ojcowskiej twarzy,\\
	Nie chciej mi bronić Twojego natchnienia;\\
	Wróć radość, którą niewinność nas darzy \\
	I racz mię w duchu utwierdzić rządzenia.\\
\end{center}

\oremus{Sekreta}{Fac nos, quǽsumus, Dómine, his munéribus offeréndis
	conveniénter aptári: quibus ipsíus venerábilis sacraménti celebrámus
	exórdium.}{Spraw, prosimy Cię, Panie, abyśmy należycie się usposobili do
	ofiarowania Tobie tych darów, przez które obchodzimy początek czcigodnego
	misterium postu.}

\medskip

\rubric{Prefacja Wielkopostna.}

\oremus{Communio (Ps 1:2,3)}{Qui meditábitur in lege Dómini die ac nocte, dabit
	fructum suum in témpore suo.}{Kto rozważa Prawo Pańskie dniem i nocą, ten wyda
	owoc w swym czasie.}

\proroctwo{Pieśń}

\gregorioscore{Melody/ABC}

\bigskip

\oremuss{
	2. Scrutátor alme córdium, \\
	Infírma tu scis vírium: \\
	Ad te revérsis éxhibe \\
	Remissiónis grátiam.\\

}{

	3. Multum quidem peccávimus,\\
	Sed parce confiténtibus:\\
	Ad nóminis laudem tui\\
	Confer medélam lánguidis.\\
}

\newpage

\oremuss{
	4. Concéde nóstrum cónteri\\
	Corpus per abstinéntiam;\\
	Culpæ ut relínquant pábulum\\
	Ieiúna corda críminum.\\

}{

	5. Præsta, beáta Trínitas,\\
	Concéde, simplex Unitas;\\
	Ut fructuósa sint tuis\\
	Ieiuniórum múnera.\\
	Amen.
}

\proroctwo{Pieśń}

\oremuss{
	1. Witam Cię, witam Przenajświętsze Ciało,\\
	któreś na krzyżu sromotnie wisiało\\
	Za nasze winy, Synu jedyny,\\
	Ojca wiecznego, Boga prawego,\\
	Zbawicielu świata.

	\medskip

	2. Dajem Ci pokłon, Bogu prawdziwemu, \\
	W tym Sakramencie dziwnie zakrytemu;\\
	Żebrzem litości i Twej miłości,\\
	Byś gniew swój srogi, o Jezu drogi,\\
	Pohamować raczył.\\
}{
	3. Zmiłuj się, zmiłuj nad nami grzesznymi,\\
	Nie racz pogardzać prośbami naszymi;\\
	Zgrom hardych siły, daj pokój miły,\\
	Wyniszcz złe rady, wykorzeń zdrady\\
	Spośród ludu swego.\\

	\medskip

	4. A kiedy przyjdzie z światem tym rozstanie,\\
	Ratuj nas, ratuj, dobrotliwy Panie!\\
	Niech z Ciała Twego, Przenajświętszego\\
	Posiłek mamy i oglądamy\\
	Ciebie łaskawego. Amen.\\
}

\bigskip

\oremus{Modlitwa po Komunii}{Percépta nobis, Dómine, præbeant sacraménta
	subsídium: ut tibi grata sint nostra ieiúnia, et nobis profíciant ad
	medélam.}{Niech nas wspomaga, Panie, przyjęty Sakrament, aby nasze posty
	były miłe Tobie, a nam przyniosły wyleczenie.}

\bigskip

\oremus{Modlitwa nad ludem}{Humiliáte cápita vestra Deo.\medskip

	Inclinántes se, Dómine, maiestáti tuæ, propitiátus inténde: ut, qui divíno
	múnere sunt refécti, coeléstibus semper nutriántur auxíliis.}{
	Pochylcie głowy wasze przed Bogiem. \medskip

	Wejrzyj łaskawie, o Panie, na wiernych, którzy chylą się przed Twoim majestatem,
	aby posileni Boskim darem zawsze byli podtrzymywani niebieską pomocą.}

\vfill

\centerline{\pgfornament[width=3cm]{106}}

\newpage

\proroctwo{Antyfona}
%
\begin{center}
	Ave, Regína cælórum,\\
	Ave, Dómina Angelórum:\\
	Salve radix, salve porta,\\
	Ex qua mundo lux est orta:\\
	Gaude, Virgo gloriósa,\\
	Super omnes speciósa,\\
	Vale, o valde decóra,\\
	Et pro nobis Christum exóra.
\end{center}

\proroctwo{Pieśń na wyjście}

\gregorioscore{Melody/AD}

\bigskip

\oremuss{
2. D\'{e}xtera P\'{a}tris, l\'{a}pis angul\'{a}ris, \\
v\'{i}a sal\'{u}tis, i\'{a}nua cael\'{e}stis, \\
\'{a}blue n\'{o}stri m\'{a}culas del\'{i}cti. \rrit \textit{Att\'{e}nde.}\\

\medskip

3. Rog\'{a}mus, D\'{e}us, t\'{u}am maiest\'{a}tem:  \\
\'{a}uribus s\'{a}cris g\'{e}mitus ex\'{a}udi:
cr\'{i}mina n\'{o}stra pl\'{a}cidus ind\'{u}lge. \rrit \textit{Att\'{e}nde.}
}{
4. T\'{i}bi fat\'{e}mur cr\'{i}mina adm\'{i}ssa:  \\
contr\'{i}to c\'{o}rde p\'{a}ndimus occ\'{u}lta: \\
t\'{u}a Red\'{e}mptor, p\'{i}etas ign\'{o}scat. \\\rrit \textit{Att\'{e}nde.}

\medskip

5. Innocens c\'{a}ptus, nec rep\'{u}gnans d\'{u}ctus, \\
t\'{e}stibus f\'{a}lsis pro \'{i}mpiis damn\'{a}tus:  \\
quos redem\'{i}sti, tu cons\'{e}rva, Chr\'{i}ste.  \\\rrit \textit{Att\'{e}nde.}
}

\vfill

\centerline{\pgfornament[width=7cm]{82}}

\vfill


\end{document}







